\chapter{Choosing a Sampler}
\label{ch:choosing-sampler}

\newthought{Seventeen methods, four families, one decision.} \Jtwo\ ships four
\emph{simple} (fixed-set) samplers, one \emph{adaptive} contour tracer, ten
\emph{\textsc{mcmc}} methods, and two \emph{nested sampling} methods. Two of the ten
\textsc{mcmc} methods --- \sampler{ToyMCMC} and \sampler{PTMCMC} --- have a full
chapter each (Chapters~\ref{ch:toymcmc}--\ref{ch:ptmcmc}); the other eight are
implemented but not yet closed for validation (\S\ref{sec:mcmc-family}). This short
chapter is the decision guide across all four families; the chapters that follow are
the per-method references for the families that have them.

\section{Sampler families and their chapters}
\label{sec:sampler-chapter-index}
\label{sec:fixed-vs-feedback}

The structural split matters more than any individual method:

\begin{description}[style=nextline, leftmargin=1.2em]
  \item[Simple methods] \mbox{}\\[-0.4\baselineskip]
        \begin{itemize}
          \item \sampler{Random} --- Chapter~\ref{ch:random}
          \item \sampler{Grid} --- Chapter~\ref{ch:grid}
          \item \sampler{CSV} --- Chapter~\ref{ch:csv-sampler}
          \item \sampler{Bridson} --- Chapter~\ref{ch:bridson}
        \end{itemize}
        These methods decide \emph{every} point before the first result returns:
        \yamlkey{point\_number} random draws, the full grid, the replayed file, or the
        blue-noise disk set (\sampler{Random}, \sampler{Grid}, \sampler{CSV},
        \sampler{Bridson}). Submission is
        embarrassingly parallel, resume is trivial (finished points are simply not
        re-proposed), and no point depends on any other.
        
  \item[The adaptive method] \mbox{}\\[-0.4\baselineskip]
        \begin{itemize}
          \item \sampler{AdaptiveBridson} --- Chapter~\ref{ch:adaptive-levelset}
        \end{itemize}
        It reacts to results toward a fixed goal: one generation waits for its
        observables at a barrier, then concentrates the next generation where the target
        level set runs (the level-set algorithm and its full reference are in that
        chapter). Fewer
        wasted evaluations, at the price of barrier synchronisation.
        
  \item[\textsc{mcmc} methods] \mbox{}\\[-0.4\baselineskip]
        \begin{itemize}
          \item Stable and fully documented: \sampler{ToyMCMC} ---
                Chapter~\ref{ch:toymcmc}; \sampler{PTMCMC} --- Chapter~\ref{ch:ptmcmc}.
          \item No dedicated chapter yet: \sampler{MCMC}, \sampler{AMMCMC}/\sampler{AM},
                \sampler{DRAM}, \sampler{EnsembleMCMC} (alias \sampler{Ensemble}),
                \sampler{DEMCMC}, and \sampler{PTEnsemble}; see
                Section~\ref{sec:mcmc-family}.
        \end{itemize}
        Together these methods explore a posterior by proposing a Metropolis step per
        chain. Every member gives you a posterior-shaped point cloud with no evidence
        estimate, at a
        fraction of nested sampling's per-point bookkeeping.
        
  \item[Nested sampling methods] \mbox{}\\[-0.4\baselineskip]
        \begin{itemize}
          \item \sampler{MultiNest} --- Chapter~\ref{ch:multinest}
          \item \sampler{Dynesty} --- Chapter~\ref{ch:dynesty}
        \end{itemize}
        These methods also react to results, but toward a different goal: they shrink a
        population of \emph{live
        points} down through the likelihood, one point at a time, until the remaining
        volume can no longer change the answer. The by-product of that bookkeeping is
        the \textbf{Bayesian evidence} ($\log Z$) alongside ordinary posterior samples
        --- something no other family here produces.
\end{description}

\section{The \textsc{mcmc} family}
\label{sec:mcmc-family}

All ten \textsc{mcmc} methods share one base implementation and one config block,
\yamlkey{Sampling.Bounds} --- the \emph{same} top-level key nested sampling uses, but
read differently. Two of them are stable, closed for validation, and fully
documented:

\begin{description}[style=nextline, leftmargin=1.2em]
  \item[\sampler{ToyMCMC}] Plain random-walk Metropolis, independent chains
        (Chapter~\ref{ch:toymcmc}) --- the family's reference method and the vehicle
        for its shared runtime.
  \item[\sampler{PTMCMC}] Parallel tempering on top of \sampler{ToyMCMC}
        (Chapter~\ref{ch:ptmcmc}): several temperature rungs run in parallel and
        periodically swap states, so hot replicas cross barriers that would trap a
        single cold chain.
\end{description}

\begin{jwarn}
The other eight \textsc{mcmc} methods below are implemented and registered --- every
one raises no \code{NotImplementedError} --- but, unlike \sampler{ToyMCMC} /
\sampler{PTMCMC} and the other three families, they do not yet have a dedicated
chapter: there is no closed \yamlkey{Bounds} validation gate or Simple/Full
\textsc{yaml} template for them yet. What follows is everything confirmed from the
source; treat it as a map of what exists, not a complete reference. The runtime
agrees with this warning: \cli{Jarvis man sampler} does not list these eight at all,
and \cli{Jarvis man sampler --full} marks every one of them \code{status: unstable}
for exactly this reason (Section~\ref{sec:cli-man}) --- \code{Jarvis validate} does
not yet reject an unknown key under their \yamlkey{Bounds}.
\end{jwarn}

These eight read the same chain-count/iteration-budget/proposal-scale keys
Chapter~\ref{ch:toymcmc} documents --- \yamlkey{Bounds.num\_chains} (alias
\yamlkey{chains}) and \yamlkey{Bounds.num\_iters} (aliases \yamlkey{iterations},
\yamlkey{steps}) set the chain count and iteration budget, \yamlkey{Bounds.proposal\_scale}
sets the Metropolis step size --- plus \yamlkey{Bounds.on\_failure} governing points
whose likelihood does not evaluate. Because their \yamlkey{Bounds} zone is not closed
yet, these alias spellings actually work at runtime; the same spellings are rejected
outright under \sampler{ToyMCMC}/\sampler{PTMCMC}'s closed zone
(Chapter~\ref{ch:toymcmc}), which only recognises the canonical names. Generation is
chain-based: one generation advances
every active chain by one step (or, for \sampler{DRAM}, one delayed-rejection stage),
waits at the usual feedback barrier for the batch's observables, and decides the next
step from there --- the same propose/barrier/absorb shape as \sampler{AdaptiveBridson},
just with Metropolis acceptance instead of level-set refinement.

\begin{description}[style=nextline, leftmargin=1.2em]
  \item[\sampler{MCMC}] Plain Metropolis --- fixed proposal scale, no adaptation;
        \sampler{ToyMCMC} (Chapter~\ref{ch:toymcmc}) is this method's documented,
        closed-validation sibling.
  \item[\sampler{AMMCMC}, \sampler{AM}] Adaptive Metropolis: the proposal covariance
        tunes itself from the chain's own history after a warm-up window.
  \item[\sampler{DRAM}] Delayed-Rejection Adaptive Metropolis: a rejected proposal gets
        a second, shrunk-scale try within the same step before moving on.
  \item[\sampler{EnsembleMCMC} (alias \sampler{Ensemble})] Complementary-half ensemble
        moves (emcee-style): each chain proposes from the spread of the other half of
        the walker population, no per-chain tuning needed.
  \item[\sampler{DEMCMC}] Differential Evolution MCMC: proposals jump along the vector
        between two other chains' current positions, also drawn from the complementary
        half of the population.
  \item[\sampler{PTEnsemble}] Parallel tempering with ensemble moves within each
        temperature rung --- \sampler{PTMCMC} and \sampler{EnsembleMCMC}'s techniques
        combined.
\end{description}

Resume is registered as \code{implemented} for these eight too --- the same bar as
\sampler{ToyMCMC}/\sampler{PTMCMC} and the simple/adaptive samplers, ahead of nested
sampling's \code{partial} status.

\section{The checklist}
\label{sec:decision-table}

Ask these in order for \emph{one concrete goal} (Section~\ref{sec:many-samplers}
explains why "one goal" and "one project" are not the same thing); the first
condition you satisfy is your answer.

\begin{enumerate}
  \item \textbf{Do you need the Bayesian evidence ($\log Z$), for a significance or a
        model comparison?} $\rightarrow$ \sampler{Dynesty} (default), or
        \sampler{MultiNest} if the card must stay byte-compatible with a V1
        \yamlkey{Method: MultiNest} card.
  \item \textbf{Are you re-running a set of points you already have} ---
        benchmarks, a paper's published points, or a previous scan's survivors?
        $\rightarrow$ \sampler{CSV}.
  \item \textbf{Do you want dense coverage concentrated on a boundary} --- a
        confidence contour, an exclusion edge --- rather than the whole space?
        $\rightarrow$ \sampler{AdaptiveBridson}.
  \item \textbf{Do you want a posterior-shaped point cloud, with no evidence
        estimate needed?} $\rightarrow$ \sampler{ToyMCMC} (Chapter~\ref{ch:toymcmc})
        for a well-behaved, plausibly unimodal posterior; \sampler{PTMCMC}
        (Chapter~\ref{ch:ptmcmc}) if it might be multimodal enough to trap a single
        chain.
  \item \textbf{Does a downstream tool need a regular, rectangular grid} to
        interpolate, contour, or Fourier-transform on? $\rightarrow$ \sampler{Grid},
        if the space is small and low-dimensional enough (Chapter~\ref{ch:grid}).
  \item \textbf{Otherwise: are you past first contact with this card} --- workflow
        and likelihood already debugged, ready for a map you will interpolate, plot,
        or train a surrogate on? $\rightarrow$ \sampler{Bridson}.
  \item \textbf{Otherwise} --- first look at a new space, debugging a new card, or
        any time the exact attempted-sample count matters more than even coverage:
        $\rightarrow$ \sampler{Random}.
\end{enumerate}

\section{One project, many samplers}
\label{sec:many-samplers}

A task card fixes one \yamlkey{Method}, but a \emph{project} rarely needs only one.
The same \yamlkey{Calculators}/\yamlkey{Operas}/\yamlkey{Likelihood} blocks serve
several different goals over a project's life, and each goal gets its own card,
walked through the checklist above on its own:

\begin{itemize}
  \item a \sampler{Random} or \sampler{Bridson} card to debug the workflow and see
        the landscape;
  \item an \sampler{AdaptiveBridson} card once you know which contour or exclusion
        edge you actually care about;
  \item a \sampler{ToyMCMC}/\sampler{PTMCMC} or \sampler{Dynesty}/\sampler{MultiNest}
        card when you need the posterior itself, or its evidence, for the paper;
  \item a \sampler{CSV} card to re-run the interesting points from any of the above
        under a heavier workflow, keeping their original \textsc{uuid}s
        (Chapter~\ref{ch:csv-sampler}).
\end{itemize}

None of this changes the workflow blocks --- only \yamlkey{Sampling} differs between
cards, so switching goals is a matter of copying the card and changing one block, not
rebuilding the project.

\section{Rules of thumb}

\begin{itemize}
  \item \textbf{Dimensions.} \sampler{Grid} dies exponentially --- $20$ points per axis
        in $5$ dimensions is $3.2$ million samples. \sampler{Bridson} is practical to
        $d\approx4$; \sampler{AdaptiveBridson} accepts $2 \le d \le 5$;
        \sampler{Random}, \sampler{CSV}, the \textsc{mcmc} family, and nested sampling
        have no intrinsic dimension limit.
  \item \textbf{Budget honesty.} Only \sampler{Random} and \sampler{CSV} let you dial an
        exact sample count. \sampler{Grid} is a product; \sampler{Bridson}'s count
        emerges from \yamlkey{radius}; \sampler{AdaptiveBridson} stops on convergence
        or its caps; the \textsc{mcmc} family runs a fixed chain count for a fixed
        iteration budget, not a target accepted-sample count. And a \yamlkey{selection}
        cut \emph{reduces} accepted counts for every fixed-set method
        (Section~\ref{sec:selection}).
  \item \textbf{The workflow does not care.} Workers run the same execution plan
        whatever proposed the point --- switching \yamlkey{Method} later never touches
        the \yamlkey{Calculators}/\yamlkey{Operas}/\yamlkey{Likelihood} blocks.
  \item \textbf{Nested sampling is the odd one out on stopping.} \sampler{Dynesty} and
        \sampler{MultiNest} do not stop at a target sample count or a converged
        contour --- they stop when their own evidence estimate has converged
        (Chapters~\ref{ch:multinest}--\ref{ch:dynesty}). Budget them by wall time,
        not by a point count you choose in advance.
  \item \textbf{Evidence is nested sampling's alone.} If you need $\log Z$ for a model
        comparison, only \sampler{Dynesty}/\sampler{MultiNest} produce it; the
        \textsc{mcmc} family gives you the posterior shape, not a normalising evidence.
\end{itemize}

\begin{jtip}
Every fixed-set, adaptive, and \textsc{mcmc} method mints deterministic \textsc{uuid}s
from \code{(method, seed, \ldots)} --- fix \yamlkey{Bounds.seed} and the run is re-proposable:
resume after a crash, replay on another machine, extend a \textsc{csv} of survivors, all
without double-computing a point (Chapter~\ref{ch:reproducibility}).
\sampler{Dynesty}/\sampler{MultiNest} are the exception --- their proposal \textsc{uuid}s
are fresh per point, not seed-deterministic (Chapters~\ref{ch:multinest} and
\ref{ch:dynesty}).
\end{jtip}
